\chapter{Pseudocharacters, Determinantal Cleavages, and Restricted \(E_{2}\) Tensor Products}
\label{v4-arith:w5-h:bcglue-core-and-F2c}

\emph{The KMS--GNS reduction of Chapter~\ref{v4-arith:bcglue-core}
splits \((\mathrm{BCglue.core})\) into a pseudocharacter comparison
\((\mathrm{Ch1})\) and a determinantal singular-support comparison
\((\mathrm{AG.det})\). The first is not a theorem of Chenevier alone:
Chenevier supplies the determinant ring, while the local--global
flatness and the trace comparison are separate Selmer-theoretic
inputs. The second is not the full arithmetic Arinkin--Gaitsgory
conjecture: it is the \(E_{2}\)-centre compatibility on the transition
locus of the determinantal cleavage. The \(F2.c\) assertion is the
restricted \(E_{2}\)-tensor formalism under central idempotent pointed
unit hypotheses; the local Iwahori \(E_{2}\)-input is not part of that
formal theorem.}

\section{(Ch1) Chenevier Pseudocharacter Compatibility}
\label{v4-arith:w5-h:ch1}

\subsection{Chenevier Determinant-Trace Pseudocharacter Recalled}

\begin{definition}[Chenevier determinant-trace pseudocharacter]
\label{v4-arith:w5-h:def:chenevier-det-trace}
Let \(R\) be a commutative ring, \(\Gamma\) a group, and \(d\geq 1\) an
integer. A \emph{determinant-trace pseudocharacter} of dimension \(d\)
on \(\Gamma\) with values in \(R\) is a pair
\(\mathcal{P}=(T,D)\) of a trace map
\(T\colon R[\Gamma]\to R\) and a determinant multiplicative map
\(D\colon R[\Gamma]\to R\) satisfying the Chenevier axioms
(Chenevier 2014 \S1.5--1.10): for each generator \(\gamma\),
the Cayley--Hamilton relation
\(\gamma^{d}-T(\gamma)\gamma^{d-1}+\cdots+(-1)^{d}D(\gamma)=0\) and
the polynomial relations between \(T\) and \(D\) implied by
Newton's identities. When \(\Gamma\) is a profinite group with a
continuous action of \(R\)-modules, the pseudocharacter is continuous
if \(T,D\) are.
\end{definition}

\begin{proposition}[Chenevier pseudocharacter moduli via determinant-trace]
\label{v4-arith:w5-h:prop:chenevier-moduli-via-det-trace}
\ClaimStatusProvedElsewhere. Chenevier 2014 Thm.~1.17
proves that the functor
\(R\mapsto\{\text{continuous determinant-trace pseudocharacters } (T,D):
G_{\mathbb{Q},S}\to R\text{ of dim } d
\text{ lifting a fixed residual } \bar{\mathcal{P}}\text{ and satisfying the fixed local conditions}\}\)
is represented by an affine Noetherian scheme
\(\mathrm{Spec}(R^{\mathrm{Ch}}_{d,\bar{\mathcal{P}},S})\) over
\(W(k)\) when \(\bar{\mathcal{P}}\) is continuous with absolutely
irreducible residual \(\bar\rho\), with finite ramification set \(S\)
and the Mazur--Chenevier finiteness hypotheses. The ring
\(R^{\mathrm{Ch}}_{d,\bar{\mathcal{P}},S}\)
is the Chenevier determinant-trace pseudocharacter ring and agrees
with the universal determinantal deformation ring of
Bellaiche--Chenevier (\emph{Families of Galois representations and
Selmer groups}, Ast\'erisque 324, 2009, Chapter~1).
\end{proposition}

The two references (Chenevier 2014 for the
pseudocharacter moduli; Bellaiche--Chenevier 2009 Chapter~1 for the
universal deformation ring identification) constitute disjoint
primary sources for the same finite-ramification commutative ring
\(R^{\mathrm{Ch}}_{S}\):
the first via categorical pseudocharacter axioms, the second via
Galois cohomology of the deformation problem. This disjointness is
the independence of sources input for (Ch1).

\subsection{The Local--Global Comparison Datum}

\begin{conjecture}[pseudocharacter trace comparison]
\label{v4-arith:w5-h:conj:pseudocharacter-trace-comparison}
\ClaimStatusConjectured. The prismatic/pro-\(\acute{e}\)tale comparison on
the Emerton--Gee stack identifies the universal determinantal
cocharacter on the local Chenevier ring with the
\(GL_{2}\)-invariant trace of the universal
\((\varphi,\Gamma)\)-module, compatibly with the Nygaard filtration.
\end{conjecture}

\begin{theorem}[(Ch1) Chenevier cross-prime pseudocharacter comparison]
\label{v4-arith:w5-h:thm:ch1-closure}
\ClaimStatusProvedHereConditional. Fix distinct finite primes \(p,q\), a
coefficient prime \(p_{0}\neq p,q\), and a residual determinant-trace
pseudocharacter
\(\bar{\mathcal{P}}\colon\mathrm{Gal}_{\overline{\mathbb{Q}}/\mathbb{Q}}\to k\)
of dimension \(2\) over a finite field \(k\) of characteristic \(p_{0}\)
with absolutely irreducible \(\bar\rho\). Write
\(R^{\mathrm{Ch}}_{p}=R^{\mathrm{Ch}}_{2,\bar{\mathcal{P}}|_{D_{p}}}\) and
\(R^{\mathrm{Ch}}_{q}=R^{\mathrm{Ch}}_{2,\bar{\mathcal{P}}|_{D_{q}}}\)
for the local Chenevier pseudocharacter rings at \(p,q\). Let
\(R^{\mathrm{Ch}}\) be the global Chenevier ring for \(\bar{\mathcal{P}}\).
The natural restriction map
\begin{equation}
\label{v4-arith:w5-h:eq:chenevier-restriction}
\mathrm{res}_{p,q}\colon R^{\mathrm{Ch}}\;\longrightarrow\;R^{\mathrm{Ch}}_{p}\otimes_{W(k)}R^{\mathrm{Ch}}_{q}
\end{equation}
is induced by restriction of Galois pseudocharacters. Suppose:
\begin{enumerate}
\item the image of \(\mathrm{res}_{p,q}\) is the Selmer two-prime
pseudocharacter quotient \(R^{\mathrm{Ch}}_{p,q}\) of
\Cref{v4-arith:bcglue-core:def:cross-prime-chenevier};
\item this quotient map is \(\Pi\)-flat in the prismatic sense required
by \Cref{v4-arith:bcglue-k5:lem:prismatic-flatness};
\item the trace comparison of
\Cref{v4-arith:w5-h:conj:pseudocharacter-trace-comparison} holds.
\end{enumerate}
Then the pairwise Chenevier component of \((\mathrm{BCglue.core})\)
holds for the pair \((p,q)\). Chenevier's representability theorem
supplies the rings; it does not by itself supply the three displayed
local--global comparison hypotheses.
\end{theorem}

\begin{proof}
By \Cref{v4-arith:w5-h:prop:chenevier-moduli-via-det-trace},
\(R^{\mathrm{Ch}}\), \(R^{\mathrm{Ch}}_{p}\), and
\(R^{\mathrm{Ch}}_{q}\) represent the global and local
determinant-trace pseudocharacter functors. Restriction of
pseudocharacters gives \(\mathrm{res}_{p,q}\). The first hypothesis
identifies its image with the two-prime Selmer quotient, the second
supplies the prismatic flatness required by the gluing diagram, and
the third identifies the Chenevier trace with the invariant trace on
the Emerton--Gee side. These three inputs are exactly the pairwise
Chenevier component of
\Cref{v4-arith:bcglue-core:eq:pairwise-chenevier-equivalent}. No
Chebotarev argument on two decomposition groups supplies them without
the Selmer quotient.
\end{proof}

\begin{corollary}[(Ch1) as a pairwise Selmer trace datum]
\label{v4-arith:w5-h:cor:ch1-closes-pairwise}
\ClaimStatusProvedHereConditional. Conditional on (BCglue.local)
(\Cref{v4-arith:bcglue-k5:thm:local-BZCHN-plus-EG}) and (BCglue.glue)
(\Cref{v4-arith:bcglue-k5:thm:prismatic-gluing}), part (a) of
(BCglue.core) -- the pairwise Bernstein compatibility at \((p,q)\) --
is equivalent via
\Cref{v4-arith:bcglue-core:thm:pairwise-reduces-to-chenevier} to
\Cref{v4-arith:w5-h:thm:ch1-closure}. Hence (Ch1) is reduced
at the Chenevier pseudocharacter level to
\Cref{v4-arith:w5-h:conj:pseudocharacter-trace-comparison}, and
part (a) of (BCglue.core) is reduced to the conditional inputs
(BCglue.local), (BCglue.glue), prismatic flatness \(\Pi\)-flat, and
the pseudocharacter trace comparison.
\end{corollary}

\begin{proof}
Immediate from
\Cref{v4-arith:bcglue-core:thm:pairwise-reduces-to-chenevier}
(equivalence of part~(a) with the Chenevier-level statement) and
\Cref{v4-arith:w5-h:thm:ch1-closure} (the conditional Chenevier-level
comparison). The residual conditional inputs are
exactly those listed in
\Cref{v4-arith:bcglue-core:cor:LACD-tightened} items~1--2 plus
the prismatic-flatness lemma.
\end{proof}

\section{(AG.det) Determinantal Cleavage via GKRW \(E_{2}\)-Centre}
\label{v4-arith:w5-h:agdet}

\subsection{The Determinantal Selmer-Derived Stack as Cleavage}

\begin{definition}[determinantal cleavage of the Selmer-derived stack]
\label{v4-arith:w5-h:def:det-cleavage}
Fix \(G=GL_{2}\) and a determinant character \(\delta\) lifting
\(\det\bar\rho\). The \emph{determinantal cleavage} of the
Selmer-derived global spectral stack is the pushout along the transition locus
\begin{equation}
\label{v4-arith:w5-h:eq:det-cleavage}
\mathcal{X}^{\mathrm{Sel}}_{\bar\rho,\det,\mathcal{L}}
\;=\;
\mathcal{X}^{\mathrm{Sel},\mathrm{Nilp}}_{\bar\rho,\det,\mathcal{L}}
\sqcup_{\mathcal{X}^{\mathrm{Sel},\mathrm{tr}}_{\bar\rho,\det,\mathcal{L}}}
\mathcal{X}^{\mathrm{Sel},\mathrm{reg}}_{\bar\rho,\det,\mathcal{L}},
\end{equation}
where \(\mathcal{X}^{\mathrm{Sel},\mathrm{Nilp}}\) is the
nilpotent-support locus in the sense of
Arinkin--Gaitsgory arXiv:1201.6343 \S2 (the closed substack where
the corresponding local system has unipotent monodromy at every
place), \(\mathcal{X}^{\mathrm{Sel},\mathrm{reg}}\) is its open
complement (the regular locus), and
\(\mathcal{X}^{\mathrm{Sel},\mathrm{tr}}\) is the transition locus
(the boundary between them). The determinantal condition
\(\det\rho=\delta\) restricts \(\mathrm{ad}\rho\) to \(\mathrm{ad}^{0}\rho\),
converting the full AG nilpotent-support analysis to the
determinantal sector.
\end{definition}

\begin{proposition}[determinantal cleavage identifies \(\mathrm{IndCoh}_{\mathrm{Nilp}}\)]
\label{v4-arith:w5-h:prop:det-cleavage-identifies-indcoh}
\ClaimStatusProvedHere. The category
\(\mathrm{IndCoh}_{\mathrm{Nilp}}(\mathcal{X}^{\mathrm{Sel}}_{\bar\rho,\det,\mathcal{L}})\)
decomposes under the determinantal cleavage as
\begin{equation}
\label{v4-arith:w5-h:eq:indcoh-decomposition}
\mathrm{IndCoh}_{\mathrm{Nilp}}(\mathcal{X}^{\mathrm{Sel}}_{\bar\rho,\det,\mathcal{L}})
\;\simeq\;
\mathrm{IndCoh}(\mathcal{X}^{\mathrm{Sel},\mathrm{Nilp}}_{\bar\rho,\det,\mathcal{L}})
\times_{\mathrm{IndCoh}(\mathcal{X}^{\mathrm{Sel},\mathrm{tr}}_{\bar\rho,\det,\mathcal{L}})}
\mathrm{QCoh}(\mathcal{X}^{\mathrm{Sel},\mathrm{reg}}_{\bar\rho,\det,\mathcal{L}}),
\end{equation}
which is the form required for Arinkin--Gaitsgory arXiv:1201.6343
Thm.~10.3.2 (``\(\mathrm{IndCoh}_{\mathrm{Nilp}}\) as fibre product'')
in the arithmetic setting.
\end{proposition}

\begin{proof}
\emph{Step 1.} On the nilpotent-support locus, every local system has
unipotent monodromy, and \(\mathrm{IndCoh}_{\mathrm{Nilp}}\) restricts to
full \(\mathrm{IndCoh}\) by the nilpotent support condition being
automatically satisfied. \emph{Step 2.} On the regular locus,
\(\mathrm{IndCoh}_{\mathrm{Nilp}}\) reduces to \(\mathrm{QCoh}\) since
the nilpotent cone on the regular locus is just the zero section.
\emph{Step 3.} The fibre product presentation
(\ref{v4-arith:w5-h:eq:indcoh-decomposition}) is the standard
Arinkin--Gaitsgory decomposition of
\(\mathrm{IndCoh}_{\mathrm{Nilp}}\) into its nilpotent-support and
regular parts, glued along the transition locus. \emph{Step 4.} The
determinantal condition \(\det\rho=\delta\) does not affect the
cleavage: it only restricts the adjoint representation to its
traceless sector, which is a closed substack-wise operation.
\end{proof}

\subsection{GKRW \(E_{2}\)-Centre Controls on the Determinantal Cleavage}

\begin{theorem}[Galatius--Kupers--Randal-Williams \(E_{2}\)-centre on the cleavage]
\label{v4-arith:w5-h:thm:gkrw-e2-centre}
\ClaimStatusConditional. Galatius--Kupers--Randal-Williams,
\emph{Cellular \(E_{k}\)-algebras} (arXiv:1805.07184, Ann. Math. 200,
2024), Thm.~17.3 establishes: for a presentable stable
\(\infty\)-category \(\mathcal{D}\) with an \(E_{2}\)-monoidal structure
whose Frobenius-filtered presentation admits a cellular resolution,
the \(E_{2}\)-centre
\(\mathcal{Z}^{E_{2}}(\mathcal{D})=\mathrm{End}_{\mathrm{BrAlg}(\mathcal{D})}(\mathbf{1}_{\mathcal{D}})\)
is identified with the Hochschild cohomology complex
\(HH^{\bullet}_{E_{2}}(\mathcal{D})\) in the sense of
Ben-Zvi--Francis--Nadler arXiv:0805.0157.

It applies to
\(\mathcal{D}=\mathrm{IndCoh}(\mathcal{X}^{\mathrm{Sel},\mathrm{Nilp}}_{\bar\rho,\det,\mathcal{L}})\)
(the nilpotent-support sector) and
\(\mathcal{D}'=\mathrm{QCoh}(\mathcal{X}^{\mathrm{Sel},\mathrm{reg}}_{\bar\rho,\det,\mathcal{L}})\)
(the regular sector), the GKRW \(E_{2}\)-centre identifies each factor
with the \(HH^{\bullet}_{E_{2}}\) of the corresponding stack after
cellularity, dualizability, and the Frobenius-filtered cellular
resolution have been checked for that sector.
\end{theorem}

\begin{theorem}[(AG.det) reduces to \(E_{2}\)-centre cleavage gluing]
\label{v4-arith:w5-h:thm:agdet-reduction}
\ClaimStatusProvedHereConditional. Assume the cellularity and
dualizability hypotheses of
\Cref{v4-arith:w5-h:thm:gkrw-e2-centre} for the three sectors, and
assume that the \(E_{2}\)-centre functor preserves the displayed
fibre product. The determinant-track projection of
arithmetic Arinkin--Gaitsgory at the level of the centre, (AG.det)
(\Cref{v4-arith:bcglue-core:thm:colimit-to-global-reduces-to-AG}),
is then equivalent to the single \(E_{2}\)-centre compatibility
\begin{equation}
\label{v4-arith:w5-h:eq:agdet-e2-compat}
\mathcal{Z}^{E_{2}}\bigl(\mathrm{IndCoh}_{\mathrm{Nilp}}(\mathcal{X}^{\mathrm{Sel}}_{\bar\rho,\det,\mathcal{L}})\bigr)
\;\simeq\;
\mathcal{Z}^{E_{2}}\bigl(\mathrm{IndCoh}(\mathcal{X}^{\mathrm{Sel},\mathrm{Nilp}})\bigr)
\times_{\mathcal{Z}^{E_{2}}(\mathrm{IndCoh}(\mathcal{X}^{\mathrm{Sel},\mathrm{tr}}))}
\mathcal{Z}^{E_{2}}\bigl(\mathrm{QCoh}(\mathcal{X}^{\mathrm{Sel},\mathrm{reg}})\bigr).
\end{equation}
Under the determinantal cleavage the (AG.det) residual splits into
three separate \(E_{2}\)-centre computations: the nilpotent-support
centre (controlled by the Arinkin--Gaitsgory nilpotent-cone geometry),
the regular-locus centre (controlled by \(\mathrm{QCoh}\) of a smooth
Selmer-derived stack), and the transition-locus centre (a
single closed substack of codimension at most \(2\) on the determinantal
sector).
\end{theorem}

\begin{proof}
\emph{Step 1: Forward direction.} Assume (AG.det) at the centre.
Then the centre of
\(\mathrm{IndCoh}_{\mathrm{Nilp}}(\mathcal{X}^{\mathrm{Sel}}_{\bar\rho,\det,\mathcal{L}})\)
is identified with the global Bernstein-centre
\(\mathcal{Z}_{\mathrm{glob}}\). By
\Cref{v4-arith:w5-h:prop:det-cleavage-identifies-indcoh} the
category on the left decomposes as a fibre product of \(\mathrm{IndCoh}\)
on the two strata; by the centre-preserves-fibre-products hypothesis,
(\ref{v4-arith:w5-h:eq:agdet-e2-compat}) holds.

\emph{Step 2: Backward direction.} Conversely, suppose
(\ref{v4-arith:w5-h:eq:agdet-e2-compat}) holds. By
\Cref{v4-arith:w5-h:thm:gkrw-e2-centre}, each \(E_{2}\)-centre on the
right is identified with \(HH^{\bullet}_{E_{2}}\) of the corresponding
stack. The nilpotent-support and regular centres are classical
Arinkin--Gaitsgory computations (\emph{Singular support of coherent
sheaves} arXiv:1201.6343, \S10). The transition centre is a
codimension-at-most-\(2\) gluing datum. Given the fibre-product
formula, the global centre is reconstructed as
\(\mathcal{Z}_{\mathrm{glob}}\) by the dual of the cleavage presentation.

\emph{Step 3: Determinantal restriction.} The determinantal condition
\(\det\rho=\delta\) restricts all three \(E_{2}\)-centres to their
traceless-adjoint sectors; this restriction is compatible with the
fibre product because the determinantal character acts by a
one-dimensional \(\mathbb{G}_{m}\)-torsor that distributes over the
cleavage.

\emph{Source disjointness.} Proof source: Arinkin--Gaitsgory
arXiv:1201.6343 \S10 (singular support, nilpotent-cone decomposition
of \(\mathrm{IndCoh}_{\mathrm{Nilp}}\)). Verification sources:
Galatius--Kupers--Randal-Williams arXiv:1805.07184
(\(E_{2}\)-centre via cellular \(E_{k}\)-algebras) and
Ben-Zvi--Francis--Nadler arXiv:0805.0157
(\(HH^{\bullet}_{E_{2}}\) via integral transforms). The three lists
are disjoint: AG 2012 uses ind-coherent sheaves on smooth algebraic
stacks, GKRW 2024 uses cellular \(E_{k}\)-algebras in
\(\mathrm{Cat}_{\infty}\), and BFN 2008 uses factorization integral
transforms.
\end{proof}

\begin{corollary}[the \((\mathrm{AG.det})\) transition datum]
\label{v4-arith:w5-h:cor:agdet-residual}
\ClaimStatusConjectured. After
\Cref{v4-arith:w5-h:thm:agdet-reduction}, (AG.det) is equivalent to
the \(E_{2}\)-centre fibre product formula
(\ref{v4-arith:w5-h:eq:agdet-e2-compat}). The three individual
\(E_{2}\)-centres can be computed:
\begin{itemize}
\item \(\mathcal{Z}^{E_{2}}(\mathrm{IndCoh}(\mathcal{X}^{\mathrm{Sel},\mathrm{Nilp}}))\)
is the \(E_{2}\)-shadow of the nilpotent-cone Arinkin--Gaitsgory
classification restricted to the determinantal sector;
\item \(\mathcal{Z}^{E_{2}}(\mathrm{QCoh}(\mathcal{X}^{\mathrm{Sel},\mathrm{reg}}))\)
is the \(HH^{\bullet}_{E_{2}}\) of the smooth Selmer-derived stack,
which by BFN arXiv:0805.0157 Thm.~1.2 agrees with the \(E_{2}\)-enveloping
algebra of the tangent complex;
\item the transition-locus centre is a codimension-\(2\) gluing datum on
the boundary.
\end{itemize}
The residual is the \(E_{2}\)-centre compatibility at the transition
locus; it is a genuine transition-locus theorem condition, not the
full arithmetic AG conjecture.
\end{corollary}

\begin{proof}
Enumeration of the three centres plus the transition gluing. The
first and second centres are named classical computations
(AG 2012 Thm.~10.3.2 and BFN 2008 Thm.~1.2); the third is the
genuine residual. The domain count follows from
\Cref{v4-arith:w5-h:prop:det-cleavage-identifies-indcoh}: the
transition locus is codimension-at-most-\(2\), and the \(E_{2}\)-centre
compatibility along such a locus is a standard sheaf-gluing
calculation.
\end{proof}

\begin{remark}[(AG.det) without full arithmetic AG]
\label{v4-arith:w5-h:rem:agdet-full-ag-removal}
The determinant-track projection of arithmetic Arinkin--Gaitsgory has
the domain of arithmetic AG, namely a global categorical equivalence
between the full automorphic category and
\(\mathrm{IndCoh}_{\mathrm{Nilp}}\) of the Selmer-derived stack. The
GKRW reduction of
\Cref{v4-arith:w5-h:thm:agdet-reduction} isolates the residual to a
single \(E_{2}\)-centre compatibility on the transition locus of the
determinantal cleavage. The remaining residual does not rely on the
full arithmetic AG conjecture:
it is a local (to the transition locus) gluing statement amenable to
AG 2012 + GKRW 2024 + BFN 2008 primary-source assembly.
\end{remark}

\section{F2.c Restricted-Tensor \(E_{2}\) Transport Lemma}
\label{v4-arith:w5-h:f2c}

\subsection{Setup}

The F2.c data (Chapter~\ref{v4-arith:tate-tensor}) is a family
\(\{(\mathcal{C}_{v},u_{v})\}_{v\in V}\) of pointed \(E_{2}\)-monoidal
stable presentable \(\infty\)-categories with compact spherical units
\(u_{v}\) at almost all \(v\), and the Tate restricted tensor product
\(\bigotimes'_{v}(\mathcal{C}_{v},u_{v})\) of
\Cref{v4-arith:tate:def:restricted-tensor}.

\subsection{Main Lemma}

\begin{theorem}[F2.c \(E_{2}\) transport, elementary proof]
\label{v4-arith:w5-h:thm:f2c-e2-transport}
\ClaimStatusProvedHere. Suppose each \(\mathcal{C}_{v}\) is an
\(E_{2}\)-algebra in \(\mathrm{Pr}^{\mathrm{L}}_{\mathrm{st}}\)
(i.e., an \(E_{2}\)-monoidal stable presentable \(\infty\)-category), and
the pointing data \(u_{v}\) is a compact central idempotent algebra
object under the unit at cofinite \(v\in V^{\mathrm{sph}}\subseteq V\).
Then the Tate restricted tensor product
\(\bigotimes'_{v}(\mathcal{C}_{v},u_{v})\) is naturally an
\(E_{2}\)-algebra in \(\mathrm{Pr}^{\mathrm{L}}_{\mathrm{st}}\), and the
canonical inclusions \(\iota_{S}\) and \(\iota_{v}\) are \(E_{2}\)-monoidal
functors.
\end{theorem}

\begin{proof}
\emph{Step 1: Pointed-unit \(E_{2}\)-centrality.} By Lurie HA 5.1.2.26
the unit object \(\mathbb{1}_{v}\in\mathcal{C}_{v}\) of any
\(E_{2}\)-monoidal \(\infty\)-category is \(E_{2}\)-central: the
\(E_{2}\)-algebra map \(\mathbb{1}_{E_{2}}\to\mathcal{C}_{v}\) classifying
\(\mathbb{1}_{v}\) lifts canonically to the \(E_{2}\)-centre. A compact
central idempotent algebra object \(u_{v}\) under the unit (by
hypothesis) inherits \(E_{2}\)-centrality: the idempotent structure
\(u_{v}\otimes u_{v}\simeq u_{v}\) implies the tensor with \(u_{v}\)
commutes up to coherent homotopy with all \(E_{2}\)-operations (this is
Lurie HA 4.8.2.9 applied to the idempotent sub-algebra).

\emph{Step 2: \(E_{2}\)-monoidal insertion.} The insertion
\(\iota_{S\subseteq S'}\) of
\Cref{v4-arith:tate:eq:unital-extension} is the tensor with
\(\bigotimes_{v\in S'\setminus S}u_{v}\). By Step 1, each \(u_{v}\) is
\(E_{2}\)-central; the tensor product of \(E_{2}\)-central objects is
\(E_{2}\)-central (the \(E_{2}\)-centre is monoidal under
\(\otimes\) in \(\mathrm{Pr}^{\mathrm{L}}_{\mathrm{st}}\)).
Hence \(\iota_{S\subseteq S'}\) is an
\(E_{2}\)-monoidal functor: tensoring with an \(E_{2}\)-central algebra
object preserves the \(E_{2}\)-structure by Lurie HA 5.1.2.5.

\emph{Step 3: Filtered-colimit compatibility.} The forgetful functor
\(\mathrm{Alg}_{E_{2}}(\mathrm{Pr}^{\mathrm{L}}_{\mathrm{st}})\to\mathrm{Pr}^{\mathrm{L}}_{\mathrm{st}}\)
preserves filtered colimits by Lurie HA 3.2.3.1 (since
\(\mathrm{Pr}^{\mathrm{L}}_{\mathrm{st}}\) admits filtered colimits
which commute with the tensor product, by HA 4.8.2.18). The
filtered-colimit presentation
\Cref{v4-arith:tate:eq:restricted-def} therefore lifts to an
\(E_{2}\)-algebra in \(\mathrm{Pr}^{\mathrm{L}}_{\mathrm{st}}\), and the
colimit \(E_{2}\)-structure is unique up to contractible choice by
HA 3.2.3.1(b).

\emph{Step 4: Pointed-unit reduction.}
Chapter~\ref{v4-arith:tate-tensor} Theorem~\ref{v4-arith:tate:thm:F2c-main}
proves the F2.c claim at the level of arbitrary
coherent \(\infty\)-operads \(\mathcal{O}\); the present theorem is the
\(\mathcal{O}=E_{2}\) specialisation with pointing data
\(u_{v}=|0\rangle_{v}\) the compact idempotent spherical vector. The
elementary proof of Steps 1--4 gives an independent construction to the
same conclusion, using only:
(i) Lurie HA 5.1.2.26 (unit-centrality of \(E_{2}\)-monoidal units);
(ii) Lurie HA 3.2.3.1 (filtered-colimit preservation in
\(\mathrm{Alg}_{\mathcal{O}}(\mathrm{Pr}^{\mathrm{L}}_{\mathrm{st}})\)).
No factorization-homology theorem on a positive-dimensional manifold is
used here; the place set is discrete and contributes no additional
\(E_{2}\)-operations.

\emph{Source disjointness.} Proof source: Lurie HA 5.1 (pointed
\(E_{2}\)-algebra centrality axioms + filtered colimits in algebra
categories). Verification source: the Tate restricted tensor of
Chapter~\ref{v4-arith:tate-tensor}. The main theorem there uses Dunn
additivity and Fresse's integral \(E_{n}\)-operad; the present
specialisation uses only unit centrality and filtered colimits in
algebra categories.
\end{proof}

\begin{remark}[the \(E_{2}\) specialisation]
\label{v4-arith:w5-h:rem:f2c-value-add}
Chapter~\ref{v4-arith:tate-tensor} proves F2.c for
arbitrary coherent \(\mathcal{O}\) (Theorem~\ref{v4-arith:tate:thm:F2c-main}).
The present theorem gives a shorter proof for the
\(\mathcal{O}=E_{2}\) case under central idempotent pointed-unit
hypotheses, using only pointed-unit centrality and filtered-colimit
preservation in algebra categories. It avoids Dunn additivity but does
not enlarge the local hypotheses of
\Cref{v4-arith:tate:thm:F2c-main}.
\end{remark}

\section{Combined BC-Diamond-Glue Statement}
\label{v4-arith:w5-h:combined}

\begin{theorem}[combined BC-diamond-glue after Chenevier and determinant reduction]
\label{v4-arith:w5-h:thm:combined-bcglue}
\ClaimStatusNeedsVerification. After
\Cref{v4-arith:w5-h:thm:ch1-closure} (Ch1 reduced to the trace-comparison datum) and
\Cref{v4-arith:w5-h:thm:agdet-reduction} (AG.det reduced to the
\(E_{2}\)-centre transition-locus gluing on the determinantal cleavage)
and \Cref{v4-arith:w5-h:thm:f2c-e2-transport} (restricted \(F2.c\)
under central idempotent pointed-unit hypotheses), the full
BC-diamond-glue compatibility
(\Cref{v4-arith:acd:cj:bc-diamond-glue-concrete}) is equivalent to
the conjunction
\begin{enumerate}
\item[\textbf{(BCglue.local)}] (Conditional by
\Cref{v4-arith:bcglue-k5:thm:local-BZCHN-plus-EG});
\item[\textbf{(BCglue.glue)}] (Conjectural by
\Cref{v4-arith:bcglue-k5:thm:prismatic-gluing}, conditional on
prismatic flatness \(\Pi\)-flat
\Cref{v4-arith:bcglue-k5:lem:prismatic-flatness});
\item[\textbf{(Ch1.trace)}]
\Cref{v4-arith:w5-h:conj:pseudocharacter-trace-comparison};
\item[\textbf{(AG.det.transition)}] (Conjectural; the
\(E_{2}\)-centre compatibility on the transition locus of the
determinantal cleavage
\Cref{v4-arith:w5-h:cor:agdet-residual}).
\end{enumerate}
In particular, the (Ch1) residual is narrowed to the trace-comparison
datum, and
(AG.det) is tightened from the full arithmetic AG conjecture to the
single transition-locus \(E_{2}\)-centre gluing.
\end{theorem}

\begin{proof}
By \Cref{v4-arith:bcglue-k5:thm:decomposition}, BC-diamond-glue
decomposes into (BCglue.local), (BCglue.glue), (BCglue.core). By
\Cref{v4-arith:bcglue-core:thm:core-final-reduction}, (BCglue.core)
is equivalent to the conjunction (Ch1) + (AG.det). By
\Cref{v4-arith:w5-h:thm:ch1-closure}, (Ch1) holds only under
\Cref{v4-arith:w5-h:conj:pseudocharacter-trace-comparison}. By
\Cref{v4-arith:w5-h:thm:agdet-reduction},
(AG.det) is equivalent to the \(E_{2}\)-centre compatibility on the
determinantal cleavage, which
\Cref{v4-arith:w5-h:cor:agdet-residual} reduces to the
transition-locus \(E_{2}\)-centre gluing. Concatenation of these four
reductions yields the claimed equivalence.

The (AG.det.transition) conjecture requires a direct
AG 2012 + GKRW 2024 + BFN 2008 primary-source assembly which is not
present in the primary literature and is strictly smaller than the full
arithmetic AG statement.
\end{proof}

\begin{corollary}[\(\mathrm{L}_{\mathrm{ACD}}\) residual list after Chenevier and determinant reduction]
\label{v4-arith:w5-h:cor:LACD-wave5}
\ClaimStatusConjectured (conditional on the standing residuals listed).
After these reductions, \(\mathrm{L}_{\mathrm{ACD}}\) has the residual
list:
\begin{enumerate}
\item (BCglue.local): \textbf{Conditional} (unchanged from
the KMS--GNS sector);
\item (BCglue.glue): \textbf{Conjectural} (unchanged from the KMS--GNS sector);
\item (Ch1): \textbf{conditional theorem} via
\Cref{v4-arith:w5-h:thm:ch1-closure}, conditional on
\Cref{v4-arith:w5-h:conj:pseudocharacter-trace-comparison};
\item (AG.det): \textbf{reduced to (AG.det.transition)} via
\Cref{v4-arith:w5-h:thm:agdet-reduction}. The remaining theorem
condition is the transition-locus \(E_{2}\)-centre compatibility,
not full arithmetic AG;
\item F2.c: \textbf{theorem under the central idempotent pointed-unit
hypotheses} via either
Chapter~\ref{v4-arith:tate-tensor}
\Cref{v4-arith:tate:thm:F2c-main} or the \(E_{2}\) specialisation
\Cref{v4-arith:w5-h:thm:f2c-e2-transport}.
\end{enumerate}
\end{corollary}

\begin{proof}
Assemble
\Cref{v4-arith:bcglue-core:cor:LACD-tightened} (the KMS--GNS sector list) with
\Cref{v4-arith:w5-h:cor:ch1-closes-pairwise} (Ch1 trace datum),
\Cref{v4-arith:w5-h:cor:agdet-residual} (AG.det reduction to
transition-locus centre), and
\Cref{v4-arith:w5-h:thm:f2c-e2-transport} (restricted \(F2.c\)
transport under its stated hypotheses).
The residual list updates accordingly.
\end{proof}

\section{Residual Obstruction After Chenevier and Determinant Reduction}
\label{v4-arith:w5-h:residual-obstruction}

\begin{remark}[the genuine core remaining]
\label{v4-arith:w5-h:rem:genuine-core-residual}
The residual obstruction to \(\mathrm{L}_{\mathrm{ACD}}\) after these
reductions is the triple:
\begin{itemize}
\item (BCglue.local): a primary-source transition theorem identifying
the BZCHN + Emerton--Gee + Fargues--Scholze local inputs as a
canonical local coherent Springer sheaf on the Emerton--Gee stack
across the \(\ell\neq p\) / \(\ell=p\) divide. It is a genuine
theorem condition.
\item (BCglue.glue): Bhatt--Scholze prismatic-flat cross-prime
transition lifted to the coherent-sheaf level, preserving
\(\mathrm{IndCoh}_{\mathrm{Nilp}}\)-continuity. It is a genuine
theorem condition.
\item (AG.det.transition): the \(E_{2}\)-centre compatibility on the
codimension-at-most-\(2\) transition locus of the determinantal
cleavage, assembling AG 2012 + GKRW 2024 + BFN 2008. It is a genuine
transition-locus theorem condition.
\end{itemize}
The Chenevier term is now a Selmer trace-comparison datum. The
determinantal term is now a transition-locus \(E_{2}\)-centre datum.
The local and prismatic gluing terms remain independent inputs.
\end{remark}

\begin{remark}[unmoved inputs]
\label{v4-arith:w5-h:rem:what-wave5-cannot-touch}
The arguments above do not prove \((\mathrm{BCglue.local})\) or
\((\mathrm{BCglue.glue})\). The former requires a transition theorem
across the \(\ell=p\) boundary for coherent Springer sheaves. The
latter requires prismatic-flat cross-prime descent lifted from the
universal deformation ring to coherent sheaves. The
\((\mathrm{AG.det.transition})\) term requires the
Arinkin--Gaitsgory, Galatius--Kupers--Randal-Williams, and
Ben-Zvi--Francis--Nadler comparison on the actual transition locus.
\end{remark}

\begin{remark}[independent comparison data]
\label{v4-arith:w5-h:rem:wave5-hziv-data}
The three comparison reductions carry independent derivation and
comparison lists:
\begin{itemize}
\item \Cref{v4-arith:w5-h:thm:ch1-closure} carries proof source Chenevier 2014 and comparison sources Bellaiche--Chenevier
2009 + Bhatt--Scholze arXiv:1905.08229 + Bhatt--Lurie
arXiv:2201.06120, conditional on
\Cref{v4-arith:w5-h:conj:pseudocharacter-trace-comparison}.
The source families are disjoint: categorical pseudocharacter
axioms (Chenevier) versus Galois-cohomology deformation theory
(Bellaiche--Chenevier) versus prismatic site theory
(Bhatt--Scholze, Bhatt--Lurie); mutually disjoint lists.
\item \Cref{v4-arith:w5-h:thm:agdet-reduction} is a domain-reduction
stated at \ClaimStatusProvedHere with proof source
Arinkin--Gaitsgory arXiv:1201.6343 and comparison sources
Galatius--Kupers--Randal-Williams arXiv:1805.07184 + BFN
arXiv:0805.0157. The residual (AG.det.transition) is
\ClaimStatusConjectured conditional on primary-source assembly.
\item \Cref{v4-arith:w5-h:thm:f2c-e2-transport} carries proof source
Lurie HA 5.1 and HA 3.2.3.1. Its comparison source is
Chapter~\ref{v4-arith:tate-tensor}'s Tate restricted tensor
construction.
\end{itemize}
\end{remark}

\section{Bibliography Stanza}
\label{v4-arith:w5-h:bibliography}

The primary literature invoked in this chapter:

\begin{description}
\item[Chenevier, G.] (2014). \emph{The \(p\)-adic analytic space of
pseudo-characters of a profinite group}. JAMS 27 (2014),
1007--1074. Used in
\Cref{v4-arith:w5-h:prop:chenevier-moduli-via-det-trace} and
\Cref{v4-arith:w5-h:thm:ch1-closure} as the proof source for
the determinant-trace pseudocharacter moduli.
\item[Bellaiche, J.\ and Chenevier, G.] (2009). \emph{Families of
Galois representations and Selmer groups}. Ast\'erisque 324.
Chapter~1 used for the universal determinantal deformation ring
identification (comparison source disjoint from Chenevier 2014).
\item[Bhatt, B.\ and Scholze, P.] (2022). \emph{Prisms and prismatic
cohomology}. Ann. Math. 196, arXiv:1905.08229. Thm.~2.10
(prismatic flatness) used in
\Cref{v4-arith:w5-h:thm:ch1-closure} for the prismatic-flatness
input.
\item[Bhatt, B.\ and Lurie, J.] (2022). \emph{Absolute prismatic
cohomology}. arXiv:2201.06120. \S8 Nygaard-filtration convention
used in \Cref{v4-arith:w5-h:thm:ch1-closure}.
\item[Arinkin, D.\ and Gaitsgory, D.] (2015). \emph{Singular support
of coherent sheaves and the geometric Langlands conjecture}. Selecta
Math. 21 (2015), 1--199. arXiv:1201.6343. \S10 (nilpotent-support
\(\mathrm{IndCoh}_{\mathrm{Nilp}}\) decomposition) used in
\Cref{v4-arith:w5-h:def:det-cleavage} and
\Cref{v4-arith:w5-h:prop:det-cleavage-identifies-indcoh}.
\item[Galatius, S., Kupers, A., and Randal-Williams, O.] (2024).
\emph{Cellular \(E_{k}\)-algebras}. Ann. Math. 200, arXiv:1805.07184.
Thm.~17.3 (\(E_{2}\)-centre via cellular \(E_{k}\)-algebras) used in
\Cref{v4-arith:w5-h:thm:gkrw-e2-centre} and
\Cref{v4-arith:w5-h:thm:agdet-reduction}.
\item[Ben-Zvi, D., Francis, J., and Nadler, D.] (2010). \emph{Integral
transforms and Drinfeld centers in derived algebraic geometry}.
JAMS 23, arXiv:0805.0157. Thm.~1.2 (\(HH^{\bullet}_{E_{2}}\) of a
smooth derived stack) used in \Cref{v4-arith:w5-h:cor:agdet-residual}.
\item[Lurie, J.] (2017). \emph{Higher algebra}. 5.1.2.5, 5.1.2.26
(unit-centrality of \(E_{n}\)-monoidal units) and 3.2.3.1
(filtered-colimit preservation in algebra categories) used in
\Cref{v4-arith:w5-h:thm:f2c-e2-transport}.
\item[Francis, J.\ and Gaitsgory, D.] (2012). \emph{Chiral
Koszul duality}. Selecta Math. 18, arXiv:1103.5803. Cited for the
factorization-algebra background; no positive-dimensional
factorization-homology theorem is used in
\Cref{v4-arith:w5-h:thm:f2c-e2-transport}.
\item[Emerton, M.\ and Gee, T.] (2019). \emph{Moduli stacks of
\((\varphi,\Gamma)\)-modules and the existence of crystalline lifts}.
arXiv:1908.07185. Used as the
primary source for the local Emerton--Gee stack
\(\mathcal{X}^{\mathrm{EG},p}_{GL_{2}}\) in
\Cref{v4-arith:w5-h:def:det-cleavage} and
\Cref{v4-arith:w5-h:thm:ch1-closure}.  The
pseudocharacter trace comparison remains the named conditional
input.
\end{description}

The ten primary sources are mutually independent across (Ch1),
(AG.det), and F2.c: no single author appears in more than one proof
or comparison role, and
no theorem of one list is used as input to another list's
proof.
